\documentclass[12pt]{article}
\usepackage[utf8]{inputenc}
\usepackage[T1]{fontenc}
\usepackage{hyperref}
\title{Umeda Language Resources for K-12 Curriculum Development}
\author{Compiled by Claude Sonnet 5 (claude-sonnet-5) for Andr\'as Kornai}
\date{}
\begin{document}
\maketitle

\section{Introduction}
Umeda (ISO 639-3: upi) is a Papuan language of Sandaun Province, Papua New Guinea, in the Border family. Wikipedia treats it as a variety of Sowanda, alongside Waina and Punda, which may be separate languages. Speaker estimates range from about 1,100 to 1,500. This search found very little material for the language. No school materials were found.

\section{Primary Educational Resources}
No primers, textbooks or school programmes for Umeda were found in the search results.

\section{Linguistic and Reference Materials}
\begin{itemize}
\item Alfred Gell, \emph{Metamorphosis of the Cassowaries: Umeda society, language and ritual} (1975), an anthropological monograph that includes language material: \url{https://www.relbib.de/Record/1608574768?lng=de}
\item Glottolog bibliographic record for Gell (1975): \url{https://glottolog.org/resource/reference/id/80539}
\item Wikipedia on Sowanda, which covers the Umeda, Waina and Punda varieties: \url{https://en.wikipedia.org/wiki/Sowanda_language}
\item Endangered Languages Project entry: \url{https://endangeredlanguages.com/pt-br/node/98418}
\end{itemize}

\section{Religious and Community Resources}
\begin{itemize}
\item Joshua Project profile of the Umeda people (demographic and social background only): \url{https://joshuaproject.net/index.php/people_groups/15708}
\end{itemize}

\section{Digital and Audio Resources}
No audio, video or corpus resources were found in the search results.

\section{Curriculum Development Context}
Little published material exists for this language. A native-language curriculum would have to start from new documentation, working with the community. Sources disagree on the classification and on whether Umeda is a language in its own right, so this should be checked before planning.
\end{document}